\nuevaReceta{Gurullos con Conejo}{4 personas}{1 h 30 min}{receta:gurullos_conejo}

    \begin{minipage}[t]{0.35\textwidth}
        \small
        \section*{Ingredientes}
        \begin{itemize}[leftmargin=*, itemsep=1pt]
            \item 1 Conejo troceado
            \item 300 g de gurullos
            \item 200 g de garbanzos cocidos
            \item 1 Cebolla
            \item 1 Pimiento verde
            \item 2 Tomates maduros
            \item 3 Dientes de ajo
            \item 1 Pimiento seco
            \item 100 ml de vino blanco
            \item 1 Hoja de laurel
            \item Pimentón dulce
            \item Azafrán o colorante
            \item Aceite de oliva
            \item Agua o caldo
            \item Sal y pimienta
        \end{itemize}
    \end{minipage}
    \hfill
    \begin{minipage}[t]{0.6\textwidth}
        \small
        \section*{Preparación}
        \begin{enumerate}[itemsep=2pt, topsep=2pt]
            \item \textbf{Dorar el conejo:} Salpimentar el conejo y dorarlo en una cazuela con aceite. Retirar y reservar.
            \item \textbf{Preparar el sofrito:} En el mismo aceite, cocinar la cebolla, los ajos y el pimiento verde picados. Añadir el tomate rallado y dejar que reduzca.
            \item \textbf{Añadir el pimiento seco:} Incorporar la carne del pimiento seco previamente hidratado y una cucharadita de pimentón, removiendo sin dejar que se queme.
            \item \textbf{Guisar el conejo:} Devolver el conejo a la cazuela, añadir el laurel y el vino blanco, y dejar que se evapore el alcohol.
            \item \textbf{Preparar el caldo:} Cubrir con agua o caldo, añadir el azafrán y cocer a fuego medio durante 35-40 minutos.
            \item \textbf{Incorporar los garbanzos:} Añadir los garbanzos cocidos y comprobar el punto de sal.
            \item \textbf{Cocer los gurullos:} Incorporar los gurullos y cocer el tiempo indicado, removiendo con frecuencia para que no se peguen. Añadir más caldo caliente si fuera necesario.
            \item \textbf{Reposar:} Apagar el fuego cuando los gurullos estén tiernos y dejar reposar 5 minutos antes de servir.
        \end{enumerate}
    \end{minipage}

    \vspace{0.5em}
    \begin{tcolorbox}[colback=orange!5, colframe=orange!50, title=Consejo, size=small]
        El guiso debe quedar meloso, no seco. Los gurullos continúan absorbiendo caldo durante el reposo, así que conviene apagar el fuego cuando todavía quede algo de líquido.
    \end{tcolorbox}
\end{receta}
